\cleardoublepage
\thispagestyle{empty}
\begin{center}


{\Large\textbf{Elektron – Theorie und Anwendungen}}\\[1.2em]
{\large Dipl.-Ing.\,(FH)\,Christian Weilharter}\\[1.2em]
\textcopyright~2025, Christian Weilharter, Traunstein\\[2em]
\end{center}
\begin{flushleft}
	\begin{tabular}{@{}l l}
	
	
		\textbf{ISBN: (Print)} & 978-3-912302-12-7 \\[0.5em]
		%	\textbf{ISBN: (E-Book)} & 978-3-912302-01-1 \\[0.5em]
		\textbf{Auflage:} & 1.~Auflage 2026 \\[0.5em]
		\textbf{Satz:} & \LaTeX \\[0.5em]
		\textbf{Verlag:} & Christian Weilharter \\[0.5em]
		\textbf{Druck:} & Amazon KDP (Print-on-Demand) \\[0.5em]

		
		\textbf{Kontakt:} & \href{mailto:info@mathandphysics.de}{info@mathandphysics.de}\\[0.5em]
		\textbf{Web:} & \href{https://www.mathandphysics.de}{www.mathandphysics.de}\\
		

	\end{tabular}
\end{flushleft}

\vspace{2em}
\noindent
Alle Rechte vorbehalten. Kein Teil dieses Buches darf ohne schriftliche Genehmigung des Autors 
in irgendeiner Form reproduziert, gespeichert oder übertragen werden, 
weder elektronisch, mechanisch, durch Fotokopien, Aufnahmen noch auf andere Weise.
\begin{center}\small Printed in Germany\end{center}

\cleardoublepage
\chapter*{Vorwort}
\addcontentsline{toc}{chapter}{Vorwort}

Das Elektron gehört zu den unscheinbarsten und zugleich folgenreichsten Objekten der modernen Physik. 
Es ist unsichtbar klein, besitzt aber eine elektrische Ladung, eine Masse, einen Spin und eine zentrale Rolle beim Aufbau der Materie. 
Ohne Elektronen gäbe es keine Atome in ihrer heutigen Form, keine chemischen Bindungen, keine elektrische Leitung, keine Halbleiter, keine Computer, keine moderne Medizintechnik und keine Quantenelektrodynamik.

Dieses Buch versteht sich als Fortsetzung und Ergänzung zu meinem Buch \emph{Photon -- Theorie und Anwendungen}. 
Während dort das Lichtquant im Mittelpunkt stand, richtet sich der Blick hier auf das Elektron: jenes Teilchen, das mit dem Photon über die elektromagnetische Wechselwirkung verbunden ist. 
Photon und Elektron bilden gewissermaßen ein physikalisches Paar: Das Photon vermittelt elektromagnetische Wechselwirkungen, das Elektron trägt elektrische Ladung und reagiert auf elektromagnetische Felder. 
Gemeinsam prägen sie einen großen Teil der modernen Physik und Technik.

Der Weg zum Elektron ist dabei nicht nur eine Geschichte der Theorie, sondern vor allem auch eine Geschichte der Experimente. 
Kathodenstrahlen, Thomsons Nachweis des Elektrons, Millikans Öltröpfchenversuch, der Photoeffekt und die Entwicklung der Atommodelle zeigen, dass das Elektron nicht einfach postuliert wurde. 
Es wurde Schritt für Schritt sichtbar gemacht -- nicht mit den Augen, sondern durch Messungen, Ablenkungen, Spuren, Energien und Wirkungen.

Besonders spannend wird das Elektron dort, wo die klassische Anschauung versagt. 
In der Quantenmechanik ist es kein kleines Kügelchen auf einer festen Bahn. 
Es zeigt Wellencharakter, erzeugt im Doppelspaltversuch Interferenzmuster und wird durch eine Wellenfunktion beschrieben. 
Gleichzeitig tritt es bei einer Messung punktförmig in Erscheinung. 
Das Elektron ist daher weder ein klassisches Teilchen noch eine gewöhnliche Welle, sondern ein Quantenobjekt. 
Gerade diese Spannung zwischen Messbarkeit und Unanschaulichkeit macht es zu einem der faszinierendsten Gegenstände der Physik.

Das Ziel dieses Buches ist es, diese Entwicklung verständlich nachzuzeichnen: von der Entdeckung des Elektrons über seine Rolle in Materie, Chemie und Atomphysik bis hin zur Quantenmechanik, zur relativistischen Beschreibung durch die Dirac-Gleichung und zur Quantenelektrodynamik. 
Dabei sollen mathematische Begriffe und physikalische Modelle nicht als trockene Formeln erscheinen, sondern als Werkzeuge, mit denen sich die Struktur der Natur immer präziser erfassen lässt.

Ein besonderer Schwerpunkt liegt auch auf den Anwendungen. 
Elektronenmikroskope, Elektronenstrahlen in Medizin und Technik, Halbleiter, Supraleitung, Spintronik und Quanteninformation zeigen, dass das Elektron nicht nur ein theoretisches Konzept ist. 
Es ist eine der Grundlagen unserer technischen Welt. 
Wer das Elektron versteht, versteht einen wesentlichen Teil der modernen Zivilisation.

Dieses Buch richtet sich an Leserinnen und Leser, die Physik nicht nur auswendig lernen, sondern wirklich verstehen wollen. 
Es setzt Interesse und Neugier voraus, aber keine vollständige fachliche Vorbildung. 
Wo mathematische Begriffe notwendig sind, werden sie schrittweise eingeführt. 
Wo Modelle an ihre Grenzen stoßen, werden diese Grenzen offen benannt. 
Denn gerade in der modernen Physik ist es wichtig zu unterscheiden, was ein Bild erklärt -- und wo es irreführend wird.

Das Elektron zeigt eindrucksvoll, dass die Wirklichkeit tiefer ist als unsere Alltagserfahrung. 
Es zwingt uns, vertraute Begriffe wie Teilchen, Bahn, Ort, Welle und Messung neu zu denken. 
Damit steht es im Zentrum einer Entwicklung, die von der klassischen Physik zur Quantenmechanik und weiter zur Quantenfeldtheorie führt.

Dieses Buch lädt dazu ein, diesen Weg Schritt für Schritt mitzugehen.








\begin{flushright}
	\textit{Dipl.-Ing.(FH) Christian Weilharter} \\
	\vspace{0.5em}
	[Traunstein, 2025]
\end{flushright}

\newpage
\noindent
\bigskip
\noindent
\textbf{Lesehinweis}

Dieses Buch folgt einem durchgehenden roten Faden.
Begriffe werden eingeführt,
an Beispielen entfaltet
und am Ende der Kapitel in ihrer Bedeutung verdichtet.
Die Boxen im Text markieren dabei zentrale Punkte
wie Definitionen, historische Einordnungen, Denkfehler oder Kernaussagen.
Sie dienen der Orientierung,
ersetzen aber nicht den Fließtext,
in dem die eigentliche Argumentation entwickelt wird.

\medskip
\noindent
Zur besseren Übersicht sind die Boxentypen farblich unterschieden:
Hinweisboxen geben Orientierung,
Historienboxen ordnen Begriffe und Entwicklungen geschichtlich ein,
Matheboxen bündeln formale Inhalte,
und Didaktikboxen heben gedankliche Schlüsselpunkte hervor.
